\scanpage{31}
% letterhead
\noindent\begin{minipage}[t]{0.57\linewidth}
\centering\sffamily
UNIVERSITE DES SCIENCES ET TECHNIQUES\\
DU LANGUEDOC\\
\rule{0.55cm}{0.4pt}\\
INSTITUT DE MATHEMATIQUES\\
Place Eugène Bataillon\\
34060 MONTPELLIER CEDEX\\
\rule{0.65cm}{0.4pt}
\end{minipage}\hfill
\begin{minipage}[t]{0.37\linewidth}
Montpellier, le
\end{minipage}

\begin{center}
\setlength{\ULdepth}{1.45\baselineskip}
\uline{COURS\quad C4 1977-78 de $\mathbb A$. GROTHENDIECK}\par\vspace{.6\baselineskip}
\end{center}

\uline{Géométrie de l'icosaèdre}

Il est étrange que l'icosaèdre, étudié depuis plus de deux mille ans dans des contextes et des optiques très divers, ne semble jamais avoir été étudié systématiquement du point de vue combinatoire, ni de celui des réalisations géométriques sur des corps (voire des anneaux) généraux. Ce ``cours'' se voudrait une investigation commune sur ce thème, en partant de l'intuition combinatoire la plus naïve, exemplifiée sur des modèles en carton (avec un coloriage convenable des arètes).

Nous aborderons donc la théorie de l'i\deleted{\unsure{e}}cosaèdre d'un point de vue combinatoire d'abord, en explicitant dans le détail le dictionnaire entre la géométrie combinatoire de l'icosaèdre orienté, et celle de l'ensemble ``orienté'' à cinq éléments (les cinq couleurs coloriant les arètes !) - ce qui donnera en même temps une motivation et une inspiration géométrique pour étudier le groupe des automorphismes commun aux deux situations, savoir le groupe alterné $\handwritten{\unsure{\mathfrak A|\mathcal A}}_5$ (qui est aussi le plus petit groupe simple non abélien). La question de construire des ``réalisations géométriques'' de l'icosaèdre combinatoire, comme ``polyèdre régulier'' sur le corps des réels, ou sur tout autre corps, sera l'occasion à la fois d'approfondir la compréhension de la structure combinatoire, et de nous familiariser avec les notions de géométrie vectorielle et affine (voire de géométrie projective) sur des corps de caractéristique quelconque.
% archive sheet number: 30
